\chapter{能否与实验对上：逐装置审视观测链}
\section{匹配规则：先把模拟量变成测量量}
温度诊断不直接等于网格上的$\operatorname{tr}\vct P/(3n)$。多普勒线宽可能包含视线流速展宽、仪器展宽和杂质物种动力学；汤姆孙散射和航天器矩也有有限采样、背景与拟合响应。模拟至少应输出带空间、时间和视线响应的合成诊断。参数匹配还应包括$B_g/B_r$、$\beta_i,\beta_e$、质量比、$L/d_i$、密度非对称性、碰撞时间、驱动及壁面损失。

本文区分原文报告的误差与本报告建议的误差预算。原文没有给RMSE、协方差或独立验证集时，不根据图像补造百分比精度；图线重合也不等于参数无拟合的预测。

\section{MAST：热点位置与交换时间受到支持，热源仍不唯一}
Tanabe等测量强导场磁环合并。电子温度和密度来自300与130通道汤姆孙散射，空间分辨率约15与10毫米；离子温度来自32弦碳离子多普勒层析。二维剖面汇合多个放电及位置，不能当成单发全域瞬时测量。典型$B_t=0.3$--$0.8\,\mathrm T$、$B_r=0.07$--$0.15\,\mathrm T$，$d_i\sim0.1\,\mathrm m$，离子回旋半径小于约0.01米。\citep{Tanabe2015MAST}

\textbf{问题与尝试：}电子是否在重联中心局地升温，而离子在下游减速处升温？测量显示中心电子热点宽约2--5厘米，温度可在小于1毫秒内从约10升至200 eV；离子升温位于出流及下游。经过约4--10毫秒的物种能量交换时间，电子和离子都可出现三峰。这与双温流体的热源分离、交换和各向异性传导解释相容。

\textbf{成立条件：}Browning的图3是直接流体比较入口，但模型初始密度约$5\times10^{18}\,\mathrm{m^{-3}}$、温度10 eV以及耗散参数，并不构成对所有后续MAST放电的逐发匹配。其keV级离子峰与较低驱动实验的幅度有差别。实验导场增加使中心电子热点更尖，所测下游离子温度变化却较小，说明约束与热损失的影响不可遗漏。\citep{L223}

\textbf{剩余缺口：}缺少同一放电的完整电场、张量压力功、热流和磁能输入；原文还指出缺乏直接磁探针约束。因而位置符合支持“下游离子能量沉积/输运”这类解释，尚不能独立在黏性、压缩、多束展宽之间作唯一选择。其最强结论是空间、时序和输运约束，非普适物种份额。
\loc{Tanabe2015MAST}{PDF2--4页图1--5，核对几何、温度剖面、时间演化和导场扫描；L223的模拟温图在PDF7页图3。}

\section{HIT-SI：真实驱动、边界与两温耦合}
Kaptanoglu等把14.5--68.5 kHz注入频率、壁面条件及三维几何纳入双温Hall计算。离子Doppler与初步Thomson结果支持部分温度趋势，作者明确尚未进行严格定量验证。模型以离子黏性$550\,\mathrm{m^2/s}$、人工粒子扩散$250\,\mathrm{m^2/s}$及有效质量比100等为输入；壁面温度约3 eV，边界密度约$0.75\times10^{19}\,\mathrm{m^{-3}}$。\citep{Kaptanoglu2020}

\textbf{已解决到的程度：}双温模型能解释为什么独立$T_e$改变电阻率、进而反馈电流和驱动，同时离子可显著更热。图3把热源、交换和壁面流出放在同一预算内，比只展示$T_i/T_e$更有解释力。

\textbf{未解决：}两代码并非完全相同几何；人工粒子扩散尚未达到不敏感极限，中性、辐射和真实电路仍未完整纳入。注入器也在供热，因此积分功率不能直接命名为重联加热效率。已有勘误\href{https://doi.org/10.1063/5.0032632}{10.1063/5.0032632}的正文仍未取得，本报告不以未经更正核对的精确温度比例作为结论。这是实质版本缺口，不能用增加描述文字消除。

\section{MRX：能量账本能约束流体，但不是独立流体验证}
MRX使用磁探针重建磁场、电流，Mach探针约束离子流，Langmuir探针约束电子温度和密度，离子光谱诊断温度。电子流常由$\vct u_e=\vct u_i-\J/(en)$反演，电势和电场的估计也依赖测量与平衡关系。离子谱在有限视线与分辨率下仍须区分流动展宽；有关诊断原始论文和源数据未在本轮重分析。\citep{F679,F674,Yoo2017Heating}

\begin{longtable}{P{39mm}P{50mm}P{59mm}}
\toprule 指标 & 对称MRX（F674） & 非对称MRX（Yoo2017Heating）\\\midrule\endfirsthead
\toprule 指标 & 对称MRX（F674） & 非对称MRX（Yoo2017Heating）\\\midrule\endhead
离子电场沉积率 & $0.71\pm0.15$ MW & $0.44\pm0.09$ MW\\
电子电场沉积率 & $0.35\pm0.06$ MW & $0.35\pm0.07$ MW\\
报告中的热能变化项 & 离子$0.27\pm0.10$ MW，电子$0.26\pm0.10$ MW & 两者均约0.30 MW；不是由电功比直接给出\\
必要解释 & 热能变化项包含储存及边界焓流重排；更换控制体会变份额 & 密度非对称、分离面输运与控制体不同；不能外推对称2:1\\
\bottomrule\end{longtable}

F679还把有序离子加速、下游温升与有限碰撞联系起来，但其数值比较为PIC，不能归为双流体模拟成功。对流体研究最直接的要求是：同一控制体同时匹配物种电功、流能、温度和边界输运。只命中2:1，而没有命中近似1:1的热能变化项，实际上可能给出错误能量路径。

两项比值也不能假设测量误差独立后简单相除：它们共用密度、电场、控制体等测量输入，相关误差需要原始数据。原文部分热损失误差很大或只给上限。因此本报告保留各功率及其误差，不补造精确份额误差条。
\loc{F674 / Yoo2017Heating / F679}{F674正式版PDF6--9页、图6--8与表I；Yoo2017图10--11在PDF15页、式(10)在14页；F679图5、9--13及式(1)--(19)。}

\section{MMS近X线：电子闭合的可分辨预测}
\subsection{2019年观测与混合对照}
Wetherton等研究2015年10月31日约07:18 UT的MMS3穿越：距X点约$7$--$10d_i$，近对称且强导场，重联发生在相撞出流之间的磁通绳内部。电子分布矩与磁场沿轨迹给出$n,B,T_{e\parallel},T_{e\perp}$。作者采用$n_\infty=7.1\,\mathrm{cm^{-3}}$、$B_\infty=55.5$ nT、$T_{e\perp\infty}=32.4$ eV及$T_{e\parallel\infty}/T_{e\perp\infty}=1.22$描述来源。\citep{Wetherton2019}

\textbf{辨别力来自曲线转折：}图2a在密度约4--12 $\mathrm{cm^{-3}}$范围内，平行温度从近等温过渡为显著增加，而垂直温度只弱变；固定CGL或固定Boltzmann极限都不能覆盖整个区间。图1还比较俘获/通过边界和分布形状，因此并非只拟合单一温度峰。

混合模拟取$B_g/B_r=1.75$、总磁场归一电子$\beta_e=0.02$、$T_i/T_e=9.1$，力自由初态，$80\times40d_i$、$2048\times1024$网格、每格400个离子粒子，以显式电阻破冻结。在$t=80/\Omega_{ci}$选取宽约$3.5d_i$的出流截面与观测比较。低质量比100的PIC在密度空洞处出现额外平行温度下降；电子流体闭合没有这一明显偏差。这里是有限电子通过时间排序的比较，不能推成所有流体优于真实质量比动理学。

\textbf{验证强度：}原文称定量符合，但没有在本文提供完整残差统计、误差协方差或混合模型网格收敛序列；截面按观测宽度选择。它比仅观察热点位置强，但仍是单事件、带匹配条件的闭合支持。离子是动理学粒子，不属于完整双流体分配检验。FPI电子分布和磁场测量提供数据基础，本文未详列全部仪器响应和校准误差，本报告不凭常规仪器指标补充成本文实际分辨率。

\subsection{与2015年10月16日十矩事件不可混为同一次验证}
TenBarge2019MMS和Closure2026使用的是另一个非对称、弱导场事件。局地十矩的电子各向异性峰幅约2.5虽近观测量级，位置在磁鞘侧，观测却在磁层侧；后继梯度闭合仍有侧别和过宽区域问题。两事件的导场、非对称性和电子磁化条件不同，不能从强导场EoS的成功直接宣称弱导场十矩通过实验。\citep{TenBarge2019MMS,Closure2026}

\section{MMS远下游：参数拟合、外部压缩和闭合失败同时存在}
Wetherton等2021年分析2016年1月21日约01:06 UT的MMS2事件，约$100d_i$下游、$B_g/B_r\sim0.7$、入流略不对称。这个距离来自前序时序分析，不能当作局地直接测距。作者为两侧入流和出流分别选择外部来源参数，使状态方程匹配温度曲线；参数见表。\citep{Wetherton2021}

\begin{center}
\begin{tabular}{lrrrr}\toprule
区域 &$n_\infty$（cm$^{-3}$）&$B_\infty$（nT）&$T_{\perp\infty}$（eV）&$T_{\parallel\infty}/T_{\perp\infty}$\\\midrule
入流1 &47&30&28&0.9\\
出流 &41&29&30&1.0\\
入流2 &40&41&30&0.9\\\bottomrule
\end{tabular}
\end{center}

这些是原文表1的模型来源参数，\textbf{不能写成三个远处热库都被独立直接测量}。图6多数区间温度随局地$n,B$变化被解释，图7有三条区域特征曲线；但若干紫色区间垂直温度显著超过原闭合预测。

作者把偏差联系到局地俘获、有限通过时间与整个区域的磁压缩，增添$\Delta\mathcal E=\mu\int\partial_tB[\vct x(t)]\dd t$。深俘获例的推断能量增强约1.35，与选定$B/B_\infty\simeq1.34$相符；同时出流区额外项通常较小，最强垂直增能位于部分入流。图9--10的增能是依据分布模型反演，不能当成直接测得的独立能量通量。

\textbf{本轮最重要的限制证据：}原文明确没有$\Delta\mathcal E$对$n,B$的形式依赖，因而不能把这一步视为已完成的通用改进状态方程；压缩起源及是否延伸至X线也未确定。已有曲线符合支持“来源与俘获调节温度”，但经调参后的符合、反演增能和无参数预测是三种不同证据。垂直温升也可能来自区域压缩，不能全部记在重联局地耗散账下。

\section{实验结论矩阵}
\begin{longtable}{P{23mm}P{46mm}P{40mm}P{33mm}}
\toprule 证据 & 已支持的量 & 未通过的强断言 & 当前程度\\\midrule\endfirsthead
\toprule 证据 & 已支持的量 & 未通过的强断言 & 当前程度\\\midrule\endhead
MAST & 温度热点、交换时序、导场趋势 & 黏性热源唯一；通用分配 & 定性空间/趋势\\
HIT-SI & 两温趋势与驱动耦合 & 严格定量温度比；仅X点供热 & 条件性定性\\
MRX & 分物种电功与有限体积预算 & 2:1不可逆热化比；纯流体已验证 & 实测预算有误差\\
MMS强导场 & 各向异性温度及密度转折 & 完整双流体两物种分配 & 事件级闭合支持\\
MMS远下游 & 分区域EoS曲线和偏离位置 & 来源无拟合；局地通用垂直热源 & 拟合约束与失败并存\\
MMS弱导场十矩 & 部分结构、峰幅及能量演化 & 正确热点侧别；热化比获证 & 存在实质不吻合\\
\bottomrule\end{longtable}
